\documentclass[11pt]{beamer} % Определяем тип документа как презентацию
% в квадратных скобках размер шрифта: 8pt, 9, 10, 11 (def), 12, 14, 17, 20.

\input{preamble} 
% В преамбуле можно сделать небольшие изменения (цвет, шрифт итд.) Если есть желание изменить больше, возможно, проще выбрать другой шаблон.

% Начнем работу с презентацией

%Для начала подготовим титульный лист и оглавление
\input{title}



\begin{document} % начнем саму презентацию

\frame{\titlepage} % сделаем титульный слайд и оглавление:

\begin{frame}
    \frametitle{Оглавление}
    \tableofcontents
\end{frame}



% Далее начнем работу со слайдами:

\section{Введение}
    \begin{frame}{Введение и План}
        Краткое описане презентации и, возможно, план...\\
        Можно выделить что-то \textbf{жирным}, а что-то \textit{курсивом}
    \end{frame}

\section{Содержание}
    \begin{frame}{Содержание 1}
        \centering
	{\large На этом слайде все пункты появляются одновременно}

        \begin{itemize}
            \item Первый

            \item Второй

            \item Третий
        \end{itemize}

        {\large На самом деле можно делать и перечисление вместо списка}
        \begin{enumerate}
            \item Первый

            \item Второй

            \item Третий
        \end{enumerate}
    \end{frame}
    \begin{frame}{Содержание 2}
        \frametitle{Пример Визуального Эффекта}
            \centering
            {\large На этом слайде уже идут визуальные эффекты}
            \smallbreak

        \begin{itemize}
            \item<1-> Этот пункт видно сразу
            \item<2-> Этот пункт видно после первого клика
            \item<3> Этот пункт видно только на третий раз
            \item<4-> Вот, третий пункт уже отсутствует
        \end{itemize}
    \end{frame}
    \begin{frame}{На самом деле проще указывать название здесь, чем в frametitle}
        Здесь уже пример \pause
        как делать этот эффект, \pause
        
        если просто хочешь текст, без itemize
    \end{frame}
    \begin{frame}{Ещё есть такие темы}
        Можно \alert{выделить} красным важное слово.

        \begin{block}{Ремарка}
            Есть такая вставка
        \end{block}

        \begin{alertblock}{Note Bene!}
            Есть <<важная>> вставка
        \end{alertblock}

        \begin{examples}
            Этим блокам тоже можно поменять цвета/шрифт в преамбулу
        \end{examples}

        % \begin{problem}
            % Здесь можно задачку записать
        % \end{problem}
        % На самом деле есть еще theorem, definition, proof, collorary, example.
    \end{frame}
    \begin{frame}{Слайд с двумя колонками (мало ли)}
        $$i_t = r^* + \pi^d + \left(1+\frac{\lambda}{\phi(\lambda^2+\alpha)}\right)(\pi^e - \pi^d) + \frac{\lambda}{\phi(\lambda^2+\alpha)} u_t + \frac{1}{\phi}(g_t-y_t^*)$$
        \begin{columns}
            \column{0.05\textwidth}
            \column{0.45\textwidth}
            Это первая колонка
            \begin{itemize}
                \item Еще что-то
                \item И Ещё
            \end{itemize}

            \column{0.5\textwidth}
            А вот это вторая колонка.\\
            Она как-то хитро алайнится
        \end{columns}
    \end{frame}
    \begin{frame}{Картинка}
        \centering
        \includegraphics[scale=0.2]{lion}
        % Все о картинках можно узнать тут https://ru.overleaf.com/learn/latex/Inserting_Images
        % В том числе как сочетать их с текстом
    \end{frame}
    \begin{frame}{Таблица}
        \begin{table}[h!]
            \centering
            \begin{tabular}{| l | c r |}
                \hline
                Текст равняется слева & Центрируется & И справа \\
                \hline
                Верхнюю строчку я  & Левую часть тоже & Да и в целом \\
                отделил чертой & $\text{formula} = \frac{1}{\sin(x)}$ & обвел все в рамку \\ [1ex]
                А тут отступ сделал & вот, кстати, формулы как & Разберетесь, 100\% \\
                \hline
            \end{tabular}
        \end{table}
        % Больше про таблицы на https://www.overleaf.com/learn/latex/Tables
    \end{frame}
    \begin{frame}{Для библиографии}
        Много чего интересного можно прочитать в книгах по использованию Латеха. Например, у Дирака \cite{dirac}. Также есть интересные статьи, например, у Эйнштейна \cite{einstein}. На еще можно посмотреть сайт \cite{knuthwebsite}, и отрывок из книги \cite{knuth-fa}.\\
        Также можно попробовать самому \href{http://www.overleaf.com}{тут}  или по ссылке: \url{http://www.overleaf.com}
        % https://www.overleaf.com/learn/latex/Hyperlinks --- ссылки
    \end{frame}

\section{Заключение}
    \begin{frame}{Заключение}
        \centering
        {\LARGE Надеюсь, этот шаблон будет вам полезен.}\\
        {\tiny Вы также копировать отдельные элементы (преамбулу, шаблон таблицы, итд) и вставлять их через input.}
    \end{frame}

\section{Источники}
    \begin{frame}{Источники}
        \printbibliography
    \end{frame}

\section*{}  
    \begin{frame}
        \textcolor{darkblueNES}{\Huge{\centerline{Спасибо за внимание!}}}
    \end{frame}
\end{document}